% ============ Solutions and hints for exercises (Part 4: §10-§12) ============
% 原书 pp.290-292; 原书对 §10 的 Ex.10.1,10.6,10.7,10.10,10.11、§11 的 Ex.11.3,11.4,11.8-11.10、
% §12 的 Ex.12.3 未给解答, 编号按原书跳号.

\section*{\S\ref{sec:10}}

\noindent\textbf{10.2.}\quad Use the previous question and Theorem~\ref{thm:7.7}.

\noindent\textbf{10.3.}\quad In the proof of Theorem~\hyperref[thm:10.4]{4} we can choose $p$ to contain $(m, y)A[y]$.

\noindent\textbf{10.4.}\quad Let $0 \subset p_1 \subset p_2$ be a strictly increasing chain of prime ideals of $R$ and let
$0 \neq b \in p_1$, $a \in p_2 - p_1$; thus $ba^{-n} \in R$ for all $n \ge 0$. Take $F = \sum u_iX^i$ to be a
root of $f^2 + af + X = 0$. Then $u_1 = -a^{-1}$, and for all $i$ we have
$u_i \in a^{-2i+1}R$, so that $bF(X) \in R[[X]]$ but $F(X) \notin R[[X]]$.

\noindent\textbf{10.5.}\quad The first part comes from Theorem~\hyperref[thm:10.1]{1}. For the second part, by \S\ref{sec:9}, Lemma~\hyperref[lem:9.1]{1}
the integral closure of $R$ in $K$ is not the whole of $K$, and therefore coincides
with $R$, so $R$ is integrally closed; on the other hand, for $x \in K - R$ we have
$R[x] = K$, hence $x^{-1} \in R[x]$, so that $x^{-1}$ is integral over $R$ and hence in $R$.
Thus $R$ is a valuation ring. If $\dim R > 1$ then there is a prime ideal $p$ of $R$
distinct from $(0)$ and from the maximal ideal, and thus $R_p$ is intermediate
between $R$ and $K$.

\noindent\textbf{10.8.}\quad Let $v\colon L^* \to G$ be the additive valuation corresponding to $S$, and choose
$x_1, \ldots, x_e \in L$ such that $v(x_1), \ldots, v(x_e)$ represent the different cosets of $G'$ in
$G$, and $y_1, \ldots, y_f \in S$ such that their images in $k$ are linearly independent
over $k'$. It follows easily from the previous question that the $ef$ elements
$x_iy_j$ are linearly independent over $K$.

\noindent\textbf{10.9.}\quad If $S \subset S_1$ then the residue field $k_1$ of $S_1$ contains a valuation ring $A \neq k_1$
such that $S$ is the composite of $S_1$ and $A$. We have $k \subset A \subset k_1$, but by the
previous question $k_1$ is an algebraic extension field of $k$, hence integral
over $A$. But $A$ is integrally closed, therefore $A = k_1$, a contradiction.

\section*{\S\ref{sec:11}}

\noindent\textbf{11.1.}\quad Let $B$ be a valuation ring of $\bar{K}$ dominating $A$, and $G$ its value group. Then
for $\alpha \in m_B$ we have $v(\alpha)/n \in G$ for all $n \ge 1$, so that $G$ has no minimal element. Also it is
easy to see that some multiple of $v(\alpha)$ belongs to $v(K^*)$, so that $G$ is
Archimedean.

\noindent\textbf{11.2.}\quad If $B$ is a valuation ring of $L$ dominating $A$ and $G$ its value group, set
$H = v(K^*)$ and $e = [G:H]$. Then $x \in G \Rightarrow e \cdot x \in H$. Hence $G$ is isomorphic to a
subgroup of $H$, and $G \cong \Z$.

\noindent\textbf{11.5.}\quad Just use Forster's theorem (5.7).

\noindent\textbf{11.6.}\quad By Ex.~\ref{ex:9.7-ex}, $A = \Z[\sqrt{10}] = \Z[X]/(X^2 - 10)$. Then $A/3A = \Z[X]/(3,
X^2 - 1) = (\Z/3\Z)[X]/(X - 1)(X + 1)$, so that $P = (3, \sqrt{10} - 1)$ is a prime
ideal of $A$. This is not principal, since if $P = (\alpha)$ with $\alpha = a + b\sqrt{10}$, then
one and sees easily that the norm $N(\alpha) = a^2 - 10b^2$ would have to be $\pm 3$,
but this is impossible since the congruence $a^2 \equiv \pm 3 \bmod 5$ has no solution.

\noindent\textbf{11.7.}\quad Let $P_1, \ldots, P_r$ be the maximal ideals of $A$; choose an element $\alpha \in P_1$ such
that $\alpha \notin P_2 \cup P_3 \cup \cdots \cup P_r$. Then $\alpha A = P_1$, and similarly each of the prime
ideals is principal. Thus by Theorem~\hyperref[thm:11.6]{6} any ideal is principal. (Of course
this also follows from Theorem~\ref{thm:5.8}.)

\section*{\S\ref{sec:12}}

\noindent\textbf{12.1.}\quad Suppose that $L$ is normal over $K$, and let $G = \Aut_K(L)$. Let $S_1$ and $S$ be
valuation rings of $L$ dominating $R$, and let $S_1, \ldots, S_l$ be the conjugates
of $S_1$ by elements of $G$, and $A = S_1 \cap \cdots \cap S_l \cap S$. If $S \neq S_i$ for any $i$ then by
Ex.~\ref{ex:10.9-ex}, there are no inclusions among $S_1, \ldots, S_l, S$, and we can apply
Theorem~\hyperref[thm:12.2]{2}. Letting $n, n_i$ be the maximal ideals of $S, S_i$, and setting
$n \cap A = p$, $n_i \cap A = p_i$ we have $p_1 \cdots p_l \not\subset p$, so that we can choose
$x \in p_1 \cap \cdots \cap p_l$ with $x \notin p$. Then $x \notin n$, but since $x \in n_i$ for all $i$, all the
conjugates of $x$ over $K$ belong to $n_1$, and the coefficients of the minimal
polynomial of $x$ over $K$ belong to $n_1 \cap K = \rad(R)$. Thus it follows easily
that $x \in n$, a contradiction.

\noindent\textbf{12.2.}\quad By Ex.~\ref{ex:10.3-ex}, $\bar{R}$ is the intersection of all valuation rings of $L$ dominating $R$.
Ex.~\ref{ex:10.9-ex} can easily be extended to the infinite case, so that the second part
follows from the first. For the first part, reduce to the finite case and use the
previous question and Theorem~\hyperref[thm:12.2]{2}.

\noindent\textbf{12.4.}\quad Let $\mathcal{P}$ be the set of height $1$ prime ideals of $A$, and for $p \in \mathcal{P}$ set $I_p = a_pA_p$.
Then $x \in \tilde{I} \iff xI^{-1} \subset A_p$ for all $p \in \mathcal{P} \iff x \in a_pA_p$ for all $p \in \mathcal{P}$. Hence $\tilde{I}$ is the
intersection of $I_p \cap A$ taken over the finitely many $p \in \mathcal{P}$ such that $I_p \neq A_p$.
